\scanpage{147}
\uline{Th} Si $A_i\to A$ \unsure{ds|dans} $L^p(\mathscr H)$, alors $\widehat A_{\unsure{i}}\to\widehat A$ \unsure{ds|dans} $\Sigma_p$
\unsure{Ceci vient|Cela résulte} facilement \unsure{du|un}

\uline{lemme} Pour \unsure{tt|tout} $\varepsilon>0$, on a $\rho>0$ et $i_0$ tel que $i>i_0$, implique

\[
\sum_{r_\nu(A_i)\leq\rho}\bigl(r_\nu\unsure{(A_i)|[A_i]}\bigr)^p\leq\varepsilon
\]

Pour \unsure{s'en|nous} convaincre, soit $\unsure{w|\omega}(t)$ la fonction convexe qui est identique à $e^{pt}$ pour $t\leq\log\rho$, et linéaire pour $t\geq\log\rho$, i.e $\unsure{-}\unsure{w|\omega}=p\rho^p(t-\log\rho)+\rho^p=\rho^p\log e\left(\frac{e^t}{\rho}\right)^{\unsure{p}}$ on a (posant $r_\nu=\unsure{r|v|\nu}_\nu(A_i)$, \deleted{\unsure{$s_\nu$ analogue}}\inserted{$s_\nu=\unsure{r|v|\nu}_\nu(|A_i|)$,})

\margin{\unsure{On n'a pas besoin} de la continuité de $A\to|A|$, si on \unsure{remplace} $|A|$ par $A^*A$\\$\unsure{>}$}

\[
\sum \unsure{w|\omega}(\log\unsure{r|v|\nu}_\nu)\leq\sum \unsure{w|\omega}(\log(s_\nu))
\]

soit

\[
\sum_{r_\nu\leq\rho}r_\nu^p+\rho^p\sum_{r_\nu\unsure{\geq|>}\rho}\log e\left(\frac{r_\nu}{\rho}\right)^p
\leq\sum_{s_\nu\leq\rho}s_\nu^p+\rho^p\sum_{s_\nu\unsure{\geq|>}\rho}\log e\left(\frac{s_\nu}{\rho}\right)^p
\]

On peut donc \uline{supposer les $A_i\geq0$}, \inserted{\unsure{(}car $|A_i|\to|A|$)} et prouver que alors, pour \unsure{tt|$\forall$|tout} $\varepsilon>0$, on a $\rho>0$ et $i_0$ tel que $i>i_0$ implique

\[
\sum_{\unsure{r_\nu(A_i)|r_\nu\lvert A_i\rvert}\leq\rho}\unsure{r_\nu(A_i)^p|r_\nu\lvert A_i\rvert^p}
+\rho^p\sum_{r_\nu(A_i)\unsure{\geq|>}\rho}\log e\left(\frac{r_\nu(A_i)}{\rho}\right)^p\unsure{<|\leq}\varepsilon
\]

On prouve que d'abord, par extraction d'une \deleted{\unsure{suite|Soit} $r_\nu\unsure{=}\,\unsure{r|v|\nu}_\nu(A_{\unsure{i'|i}})$, il suffit de montrer séparément}

\margin{\textcircled{\scriptsize \unsure{si Si|$s_iS_i$}}}

\[
\deleted{\text{a)}\qquad\rho^p\sum_{r_\nu\unsure{\geq|>}\rho}\log e\left(\frac{r_\nu}{\rho}\right)^p\to0\qquad\text{si }\rho\to0}
\]

\deleted{b) \unsure{l'énoncé du lemme|l'inégalité du lemme} pour les hermitiens positifs} $\unsure{\deleted{A_i}|A_i}$.

\deleted{\unsure{Or} \illegible{7}}

\uline{Preuve de a)} Or, on a en effet

\margin{\deleted{\illegible{1}}}

\[
\begin{aligned}
\rho^{\unsure{\nu|p}}\sum_{r_\nu>\rho}\log e\left(\frac{r_\nu}{\rho}\right)^p
&=\rho^{\unsure{\nu|p}}\int_{r>\rho}^{+\infty}-\log e\left(\frac{\unsure{r_\nu|r}}{\rho}\right)^p\,dn(r)\\
&=\rho^p\left[-n(r)\log-\right]_{\rho+0}^{+\infty}
+p\rho^p\int_\rho^{+\infty}\frac{n(r)}{r}\,dr\\
&=\rho^p n(\rho)+p\rho^p\int_\rho^{+\infty}\frac{n(r)}{r}\,dr
\end{aligned}
\]

et il est bien connu, puisque $\sum r_\nu^p<\infty$, que $\rho^p n(\rho)\to0$
